% !PRINTED 668
%% [running head, not body text -- commented out 2026-08-07] Page furniture (running
%% head, folio, signature) is excluded from this edition by standing policy.
%% 668 \hfill \gesperrt{Twenty-second Chapter.} \hfill \S.\,184.
\weberpage{668}
%<<CONT
where $\varrho_n$ is a root of unity of degree $m^2$, and $\alpha_n$ is a
number in $\Omega_m$, and it would still remain to be proved that
$e(r^n)$ is the \gesperrt{square of a unit} in $\Omega_m$.

Since $e(r)$ is a real unit, we have

\weq{2}{e(r^n) = e(r^{-n}),}

and we shall further stipulate that the sign of the root is so
determined (which, upon a suitable assumption about $\varrho_n$, we
may always take to be the case) that

\weq{3}{\sqrt{e(r^n)} = \sqrt{e(r^{-n})}.}

By hypothesis the product $(r^n, x_0)\,(r^{-n}, x_0)$ is a number of
the field $\Omega_m$ which is not altered by the substitution
$(r, r^{-1})$, i.e. a real number. Further, by \S.\,174, (5),

\[
(r^{m_2 n}, \eta)\,(r^{-m_2 n}, \eta) = \pm p,
\]

and hence likewise real. In the same way $\alpha_n \alpha_{-n}$ is
real, and from this it follows by (1) that

\[
\varrho_n \varrho_{-n} = \frac{(r^n, x_0)\,(r^{-n}, x_0)}
{e(r^n)\ \alpha_n \alpha_{-n}\ \Pi(r^{m_2 n}, \eta)\,(r^{-m_2 n}, \eta)}
\]

is a real number, which, since it is a root of unity, can only be
$= \pm 1$.

Now in the equation

\[
(r, x_0)\,(r^{-1}, x_0) = \pm\ e(r)\ \alpha_1 \alpha_{-1}\ \Pi(\pm p)
\]

which results from this for $n = 1$, we can carry out every
substitution $(r, r^n)$, and from this it appears that the sign of
$\varrho_n \varrho_{-n}$ is independent of $n$, so that

\weq{4}{\varrho_n \varrho_{-n} = \varrho_1 \varrho_{-1} = \pm 1}

holds.

Now we derive further from (1):

\weq{5}{\varrho_n \varrho_1^{-1} \sqrt{e(r^n)}\ \sqrt{e(r)}^{-n}
= \frac{(r^n, x_0)\,(r, x_0)^{-n}}
{\alpha_n \alpha_1^{-1}\ \Pi(r^{m_2 n}, \eta)\,(r^{m_2}, \eta)^{-n}},}
% [source sic: (5) prints both \varrho_1^{-1} and \alpha_1^{-1}; formation from (1) and the displayed -n powers requires \varrho_1^{-n} and \alpha_1^{-n}]

and the right side is a number of the field $\Omega_m$. The left
side shows, however, that it is a unit, and accordingly, if
$\mathfrak{E}(r)$ denotes a real unit, we may set

\weq{6}{\varrho_n \varrho_1^{-1} \sqrt{e(r^n)}\ \sqrt{e(r)}^{-n} = r^k \mathfrak{E}(r).}
% [source sic: (6) repeats the inherited \varrho_1^{-1}; the derivation requires exponent -n]

If we substitute this value in formula (5) and make the substitution
$(r, r^{-1})$, there results

\weq{7}{\varrho_{-n} \varrho_{-1}^{-1} \sqrt{e(r^n)}\ \sqrt{e(r)}^{-n} = r^{-k} \mathfrak{E}(r),}
% [source sic: (7) prints \varrho_{-1}^{-1}; the substituted form of the preceding relation requires exponent -n]
%CONT>>
